\documentclass{beamer}
\usepackage{color}
\usepackage{bm}
\usepackage{graphicx}
\usepackage{hyperref}
\usepackage{listings}
\beamertemplatenavigationsymbolsempty
\definecolor{mygreen}{rgb}{0,0.6,0}
\definecolor{mymauve}{rgb}{0.58,0,0.82}
\lstset{ %
  backgroundcolor=\color{white},   % choose the background color; you must add \usepackage{color} or \usepackage{xcolor}
  basicstyle=\footnotesize,        % the size of the fonts that are used for the code
  commentstyle=\color{mygreen},    % comment style
  deletekeywords={...},            % if you want to delete keywords from the given language
  extendedchars=true,              % lets you use non-ASCII characters; for 8-bits encodings only, does not work with UTF-8
  frame=single,                    % adds a frame around the code
  keywordstyle=\color{blue},       % keyword style
  language=bash,                   % the language of the code
  rulecolor=\color{black},         % if not set, the frame-color may be changed on line-breaks within not-black text (e.g. comments (green here))
  showspaces=false,                % show spaces everywhere adding particular underscores; it overrides 'showstringspaces'
  stringstyle=\color{mymauve},     % string literal style
  title=\lstname                   % show the filename of files included with \lstinputlisting; also try caption instead of title
}
\hypersetup{
    colorlinks=true,
    urlcolor=blue
}

\graphicspath{ {./img/} {./charts/} }


\title{Custom Template Filters for a Better World}
\author{Adam Johnson - me@adamj.eu}
\date{25th March 2015}

\begin{document}


\maketitle


\begin{frame}\frametitle{Me}

    \begin{itemize}
        \item DevOps/DBA at YPlan
    \end{itemize}

\end{frame}


\begin{frame}\frametitle{Problem: Anonymous Instances}

    \begin{center}
        \includegraphics[width=10cm]{anonymous}
    \end{center}

\end{frame}


\begin{frame}\frametitle{Sneaky Amazon?}

    \begin{center}
        \includegraphics[width=10cm]{look-i-made-a-graph}
    \end{center}

\end{frame}


\begin{frame}\frametitle{Naughty Jenkins!}

    \begin{center}
        \includegraphics[width=10cm]{jeeves}
    \end{center}

\end{frame}


\begin{frame}[fragile]\frametitle{Solution}

    \begin{itemize}
        \item \textbf{clean\_resources} playbook every 15 minutes
    \end{itemize}

    \begin{lstlisting}
- name: delete anonymous instances
  ec2_sql:
    sql: DELETE FROM ec2_instances
         WHERE length(tags) = 0 AND age > 30 minutes
    \end{lstlisting}

\end{frame}


\begin{frame}[fragile]\frametitle{Step 0}

    \begin{itemize}
        \item Add the actual play
    \end{itemize}

    \begin{lstlisting}
- name: delete anonymous instances
  hosts: ec2
  tasks:
    - local_action: ec2 state=absent
                        instance_ids={{ ec2_id }}
    \end{lstlisting}

\end{frame}


\begin{frame}[fragile]\frametitle{Step 1 - WHERE length(tags) = 0}

    \begin{itemize}
        \item Host variables:
    \end{itemize}

    \begin{lstlisting}[language=C]
"ec2-54-75-123-123.eu-west-1.compute.amazonaws.com": {
    "ec2_architecture": "x86_64",
    "ec2_client_token": "",
    ...
    "ec2_tag_Name": "my_fancy_machine",
    "ec2_tag_ansible_role": "web",
    ...
}
    \end{lstlisting}

\end{frame}


\begin{frame}[fragile]\frametitle{Step 1 - WHERE length(tags) = 0}

    \begin{itemize}
        \item Drafted python:
    \end{itemize}

    \begin{lstlisting}[language=Python]
In [1]: def filter_prefix(items, prefix):
   ...:     return [x for x in items
                    if x.startswith(prefix)]

In [2]: var_names = ['ec2_tag_a', 'ec2_time',
   ...:              'ec2_tag_b', '']

In [3]: filter_prefix(var_names, 'ec2_tag_')
Out[3]: ['ec2_tag_a', 'ec2_tag_b']
    \end{lstlisting}

\end{frame}


\begin{frame}[fragile]\frametitle{Step 1 - WHERE length(tags) = 0}

    \begin{itemize}
        \item \textbf{filter\_plugins/my\_plugins.py} for filter module:
    \end{itemize}

    \begin{lstlisting}[language=Python]
def filter_prefix(items, prefix):
    return [x for x in items if x.startswith(prefix)]

class FilterModule(object):
    def filters(self):
        return {
            'filter_prefix': filter_prefix,
        }
    \end{lstlisting}

\end{frame}


\begin{frame}[fragile]\frametitle{Step 1 - WHERE length(tags) = 0}

    \begin{itemize}
        \item Add in to play:
    \end{itemize}

    \begin{lstlisting}
- name: delete anonymous instances
  hosts: ec2
  tasks:
    - when: >
        hostvars[inventory_hostname].keys()
          | filter_prefix('ec2_tag_')
          | length == 0

      local_action: ec2 state=absent
                        instance_ids={{ ec2_id }}
    \end{lstlisting}

\end{frame}


\begin{frame}[fragile]\frametitle{Step 2 - AND age $>$ 30 minutes}

    \begin{itemize}
        \item Host variables:
    \end{itemize}

    \begin{lstlisting}[language=C]
"ec2-54-75-123-123.eu-west-1.compute.amazonaws.com": {
    ...
    "ec2_launch_time": "2014-11-15T11:25:57.000Z",
    ...
}
    \end{lstlisting}

\end{frame}


\begin{frame}[fragile]\frametitle{Step 2 - AND age $>$ 30 minutes}

    \begin{itemize}
        \item Coding by Stack Overflow:
    \end{itemize}

    \begin{lstlisting}[language=Python]
def aws_age_seconds(ec2_launch_time):
    # Strip trailing subsecond part
    launch_time = ec2_launch_time[:-len('.000Z')]
    # Turn into datetime
    time_format = "%Y-%m-%dT%H:%M:%S"
    time_tuple = strptime(launch_time, time_format)
    time_dt = datetime(*time_tuple[:6])
    # Return difference in seconds
    diff = datetime.utcnow() - time_dt
    return seconds_diff.total_seconds()
    \end{lstlisting}

\end{frame}


\begin{frame}[fragile]\frametitle{Step 2 - AND age $>$ 30 minutes}

    \begin{itemize}
        \item Full play:
    \end{itemize}

    \begin{lstlisting}
- name: delete anonymous instances
  hosts: ec2
  tasks:
    - when: >
        ec2_launch_time|aws_age_seconds > 3600 and
        hostvars[inventory_hostname].keys()
          | filter_prefix('ec2_tag_')
          | length == 0
      local_action: ec2 state=absent
                        instance_ids={{ ec2_id }}
    \end{lstlisting}

\end{frame}


\begin{frame}\frametitle{Problem: \textbf{SOLVED}}

    \begin{center}
        \includegraphics[width=10cm]{sad-anonymous}
    \end{center}

\end{frame}


\begin{frame}\frametitle{Thank you}

    \begin{itemize}
        \item \url{me@adamj.eu}
    \end{itemize}

\end{frame}


\end{document}
